% Generated by scripts/assemble.py; corrections belong in source/corrections.json.
\examhead{162}{广西大学 2024 年数学分析试卷(回忆版)}{154}
% p155-b17; source PDF page 155
\blockanchor{p155-b17}
\question{1}求极限 $\lim_{x\to 0}\frac{\sqrt{1 + \tan x} - \sqrt{1 - \sin x}}{x}$.\par
% p155-b18; source PDF page 155
\blockanchor{p155-b18}
\question{2}设 $z = x^{2} + y^{2}$, 其中 $y = f(x)$ 是由方程 $x^{2} - xy + y^{2} = 1$ 确定的隐函数, 求 $\frac{\mathrm{d}^2z}{\mathrm{d}x^2}$.\par
% p155-b19; source PDF page 155
\blockanchor{p155-b19}
\question{3}未知.\par
% p155-b20; source PDF page 155
\blockanchor{p155-b20}
\question{4}抛物线 $y^{2} = 2x$ 将 $x^{2} + y^{2} \leq 8$ 分成左右两部分, 求左右两部分面积之比.\par
% p155-b21; source PDF page 155
\blockanchor{p155-b21}
\question{5}求 $y^{2} = 2px$ 上一点, 使得此点处的法线被抛物线截取的线段最短.\par
% p155-b22; source PDF page 155
\blockanchor{p155-b22}
\question{6}求曲线积分 $\int_{L} xy \, \mathrm{d}s$, 其中 $L$ 为椭圆 $\frac{x^2}{a^2} + \frac{y^2}{b^2} = 1$ 在第一象限的部分.\par
% p155-b23; source PDF page 155
\blockanchor{p155-b23}
\question{7}求重积分 $\iiint_{V} (x^{2} + y^{2} + z) \, \mathrm{d}x \, \mathrm{d}y \, \mathrm{d}z$, 其中 $V$ 是直线 $z = y, x = 0$ 绕 $z$ 轴旋转所得曲面与 $z = 1$ 围成的区域.\par
% p155-b24; source PDF page 155
\blockanchor{p155-b24}
\question{8}求  $y = \arcsin(\cos x)$  的傅里叶级数展开式.\par
% p156-b00; source PDF page 156
\blockanchor{p156-b00}
\question{9}求曲面积分 $\iint_{S} (x + y + z) \, \mathrm{d}S$, 其中 $S: x^{2} + y^{2} + z^{2} = a^{2}, z \geq 0$.\par
% p156-b01; source PDF page 156
\blockanchor{p156-b01}
\question{10}已知 $\lim_{n\to \infty}a_n = a$ ，若 $a_{n},a > 0$ ，证明： $\lim_{n\to \infty}\sqrt[n]{a_n} = 1.$\par
% p156-b02; source PDF page 156
\blockanchor{p156-b02}
\question{11}设 $f(x)$ 在 $[a, +\infty)$ 上连续，且 $\lim_{x \to +\infty} f(x)$ 存在，证明：$f(x)$ 在 $[a, +\infty)$ 上一致连续.\par
% p156-b03; source PDF page 156
\blockanchor{p156-b03}
\question{12}证明 $f(x) = x^3 e^{-x^2}$ 是有界函数.\par
% p156-b04; source PDF page 156
\blockanchor{p156-b04}
\question{13}证明 $\int_0^{+\infty} \frac{\sin xy}{y} \, \mathrm{d}y$ 在 $[\delta, +\infty) (\delta > 0)$ 上一致收敛，在 $(0, +\infty)$ 上不一致收敛.\par
% p156-b05; source PDF page 156
\blockanchor{p156-b05}
\question{14}由闭域套定理证明 $\mathbb{R}^2$ 中的有界无限点集 $E$ 在 $\mathbb{R}^2$ 中至少有一个聚点.\par
% p156-b06; source PDF page 156
\par\Needspace{6\baselineskip}
\blockanchor{p156-b06}
\question{15}设 $S$ 是包围 $V$ 的光滑曲面, 证明 $V$ 的体积 $\Delta V$ 为\par
\nopagebreak[4]
% p156-b07; source PDF page 156
\blockanchor{p156-b07}
\displayformula{\Delta V = \frac {1}{3} \iint_ {S} (x \cos \alpha + y \cos \beta + z \cos \gamma) d S.}
% p156-b08; source PDF page 156
\blockanchor{p156-b08}
其中  $\cos\alpha,\cos\beta,\cos\gamma$  是曲面 S 外法线的方向余弦.\par
